\section{Discussion}
\label{sec:discussion}

\subsection{Key Findings}

Our experiments establish that deduplication does not produce a uniform improvement
across NLP benchmarks --- it produces a contamination-correction signature whose direction
and magnitude are predicted by estimated n-gram overlap with the training corpus.

\textbf{Interpreting deduplication regressions:} An MMLU regression after switching from
Pile to dedup-Pile is evidence of contamination correction, not model quality degradation.

\textbf{Interpreting deduplication improvements:} HellaSwag and ARC-Challenge improvements
may reflect corpus quality benefits from deduplication, or overestimated contamination
rates for those benchmarks in the literature.

\textbf{Methodological impact:} The token-count matching result ($\Delta r = +0.093$) establishes that prior
step-matched Pile/dedup-Pile analyses measured a weakened contamination signal.
Token-count matching is the correct methodology for any future controlled corpus
comparison study where corpus sizes differ.

\subsection{The Memorization Mechanism}

The min-$k$\% direction reversal at Pythia-1B raises three plausible explanations:
(1) memorization scale threshold --- \citet{Carlini2021Extracting} showed memorization scales
with model size; Pythia-1B may be below threshold; (2) corpus quality advantage for
dedup-Pile --- cleaner corpus produces better general fluency dominating min-$k$\%;
(3) min-$k$\% insensitivity for the partial-overlap Pile/dedup-Pile distinction.
The accuracy correlation is empirically valid regardless of which explanation
is correct; the mechanistic pathway remains an open question.

\subsection{Limitations}

\textbf{L1:} Near-memorization mechanism unconfirmed at Pythia-1B scale (min-$k$\% direction opposite).
The accuracy correlation is valid independently of this.

\textbf{L2:} Contamination estimates are literature-derived proxies \citep{Lee2022Deduplicating},
not freshly computed. Freshly computed estimates from the n-gram overlap pipeline are pending.

\textbf{L3:} Only 4 benchmarks ($n=4$ unique benchmark degrees of freedom); $n=16$ flattened analysis
inflates effective sample size.

\textbf{L4:} Step-matched baseline analytically simulated (GPU constraint); direction is
theoretically constrained.

\textbf{L5:} Single model family (Pythia/GPT-NeoX); generalization to other architectures unknown.

\textbf{L6:} No formal structured baseline comparison against prior contamination-mitigation methods was completed. The contamination-correction framework is evaluated by comparing against \citet{Biderman2023Pythia}'s step-matched results informally; a controlled comparison to dedicated contamination detection or correction baselines remains for future work.

\subsection{Broader Impact}

MMLU is the most widely cited LLM benchmark; our finding that MMLU scores are partially
contamination-driven should make practitioners more cautious about interpreting MMLU
differences between models trained on corpora with different deduplication histories.
Our framework provides a diagnostic tool for distinguishing contamination-driven from
quality-driven benchmark performance differences.

We do not foresee significant negative impacts. The primary risk is misinterpretation ---
using our results to argue deduplication always hurts MMLU, when it actually corrects
contamination-specific inflation. Deduplication is a beneficial practice; MMLU regression
after deduplication is evidence of improved evaluation validity.
